\documentclass[11pt,a4paper]{article}
\usepackage[margin=25mm,headheight=15pt]{geometry}
\usepackage[T1]{fontenc}
\usepackage[utf8]{inputenc}
\usepackage{lmodern,amsmath,amssymb,amsthm,mathtools,bm}
\usepackage{xcolor,graphicx,booktabs,tabularx,array,longtable}
\usepackage{microtype,enumitem,fancyhdr}
\usepackage{xurl}
\usepackage{hyperref}
\definecolor{accent}{HTML}{163B55}
\hypersetup{colorlinks=true,linkcolor=accent,citecolor=accent,urlcolor=accent,
 pdftitle={Lipschitz affine approximation forces exact martingale cotype},pdfauthor={Matthew J. Colbrook},
 pdfsubject={AIM 905: mathematical result and proof}}
\urlstyle{same}
\setlength{\parindent}{0pt}
\setlength{\parskip}{0.35em}
\setlength{\emergencystretch}{3em}
\setlist{itemsep=0.25em,topsep=0.3em}
\newtheorem{theorem}{Theorem}[section]
\newtheorem{lemma}[theorem]{Lemma}
\newtheorem{proposition}[theorem]{Proposition}
\newtheorem{corollary}[theorem]{Corollary}
\theoremstyle{remark}
\newtheorem{remark}[theorem]{Remark}
\newtheorem{definition}[theorem]{Definition}
\DeclareMathOperator{\Tr}{Tr}
\DeclareMathOperator{\supp}{supp}
\DeclareMathOperator{\Lip}{Lip}
\DeclareMathOperator{\Rea}{Re}
\newcommand{\R}{\mathbb R}
\newcommand{\C}{\mathbb C}
\newcommand{\E}{\mathbb E}
\newcommand{\N}{\mathbb N}
\newcommand{\ip}[2]{\langle #1,#2\rangle}
\pagestyle{fancy}
\fancyhf{}
\fancyhead[L]{\small\sffamily AIM 905}
\fancyhead[R]{\small\sffamily Matthew J. Colbrook}
\fancyfoot[C]{\thepage}
\renewcommand{\headrulewidth}{0.3pt}
\allowdisplaybreaks[2]
\title{\textbf{Lipschitz affine approximation forces exact martingale cotype}}
\author{Matthew J. Colbrook\\[3pt]
{\small Department of Applied Mathematics and Theoretical Physics}\\
{\small University of Cambridge, United Kingdom}\\
{\small\href{mailto:mjc249@maths.cam.ac.uk}{\texttt{mjc249@maths.cam.ac.uk}}}}
\date{3 November 2026}
\begin{document}
\maketitle
\begin{abstract}
We prove that the stated support-volume Lipschitz affine-approximation estimate forces martingale cotype at its exact exponent. The proof encodes finite dyadic martingales at separated spatial scales, constructs a power-convex equivalent norm and recovers the full martingale estimate.
\end{abstract}
\noindent\textbf{Repository target.}
This manuscript concerns AIM problem 905 at revision
\texttt{5d7e26e8444e}. The precise scope is stated below.
\medskip

\section{Main result and proof}
\label{sec:905}

The question asks whether the Lipschitz affine approximation estimate in
\cite[Question~10]{HN2019} determines the precise martingale cotype exponent.
The Sobolev version proved in \cite{HM2019} has a stronger hypothesis;
the argument below uses the stated Lipschitz estimate with support volume.

Let $X$ be a real normed space of dimension $n\geq1$, with a fixed Lebesgue
measure, and let $Y$ be a Banach space.  Write $B_X$ for its closed unit ball.
For $2\leq q<\infty$ set
\[
 \mathcal D_X(f)=\int_X\int_0^\infty
 \inf_{\Lambda\text{ affine}}
 \frac{1}{|tB_X|}\int_{x+tB_X}\|f(y)-\Lambda(y)\|^q\,dy\,
 \frac{dt}{t^{q+1}}\,dx.
\]
Here affine maps are allowed to take values anywhere in $Y$.

\begin{theorem}\label{thm:905}
Suppose there is $C<\infty$ such that every compactly supported
$1$-Lipschitz map $f:X\to Y$ satisfies
\begin{equation}\label{eq:905-hyp}
 \mathcal D_X(f)\leq C|\operatorname{supp} f|.
\end{equation}
Then $Y$ admits an equivalent uniformly convex norm whose modulus of
convexity is bounded below by a positive multiple of $\varepsilon^q$.
Consequently $Y$ has martingale cotype $q$, with the exact exponent in
\eqref{eq:905-hyp}.
\end{theorem}

For a compactly supported curve $F:\mathbb R\to Y$ define
\[
 \beta_F(s,u)^q=\inf_{a,c\in Y}\frac1{2u}
 \int_{s-u}^{s+u}\|F(v)-av-c\|^q\,dv,
 \qquad
 \mathcal D_b(F)=\int_{\mathbb R}\int_0^b
 \beta_F(s,u)^q\,\frac{du}{u^{q+1}}\,ds.
\]

\begin{lemma}[Reduction to curves]\label{lem:905-curves}
Under \eqref{eq:905-hyp}, there are $b>0$ and $K<\infty$ such that
$\mathcal D_b(F)\leq K$ for every $1$-Lipschitz curve $F$ supported in
$[0,1]$ and satisfying $\|F\|_\infty\leq1$.
\end{lemma}
\begin{proof}
Choose coordinates identifying $X$ with $\mathbb R^n$ and its chosen
Lebesgue measure with ordinary Lebesgue measure, and numbers $b,B>0$
such that
\[
 b[-1,1]^n\subseteq B_X\subseteq B[-1,1]^n.
\]
Set $\kappa=(2b)^n/|B_X|$ and
$L=\sum_{j=1}^n\|\ell_j\|_{X^*}$, where $\ell_j$ is the $j$th coordinate
functional.  Fix $R\geq\max\{1,2B\}$ and a scalar $1$-Lipschitz cutoff
$\eta$ taking values in $[0,1]$, equal to $1$ on $[-1,1]$, and zero
outside $[-2,2]$.  Define
\[
 \chi_R(x_2,\ldots,x_n)=\prod_{j=2}^n\eta(x_j/R),
 \qquad f_R(x)=F(x_1)\chi_R(x_2,\ldots,x_n).
\]
For $n=1$ the empty product equals one.  A telescoping product estimate
gives $\operatorname{Lip}_X(f_R)\leq L$ and
$|\operatorname{supp}f_R|\leq(4R)^{n-1}$.  Scaling
\eqref{eq:905-hyp} therefore yields
\begin{equation}\label{eq:905-upper}
 \mathcal D_X(f_R)\leq CL^q(4R)^{n-1}.
\end{equation}
If $x=(s,z')$, $z'\in[-R/2,R/2]^{n-1}$ and $0<t\leq1$, the cutoff
equals one throughout $x+tB_X$.  Restrict its local error integral to
the cube $x+bt[-1,1]^n$.  Averaging an affine map in the transverse
coordinates produces an affine function of its first coordinate.
Jensen's inequality consequently gives
\[
 \inf_\Lambda\frac1{|tB_X|}\int_{x+tB_X}\|f_R-\Lambda\|^q
 \geq\kappa\,\beta_F(s,bt)^q.
\]
Integrating in $s,z',t$ and substituting $u=bt$ shows that
\[
 \mathcal D_X(f_R)\geq\kappa R^{n-1}b^q\mathcal D_b(F).
\]
Together with \eqref{eq:905-upper}, this proves the lemma with
$K=CL^q4^{n-1}/(\kappa b^q)$.
\end{proof}

\begin{lemma}[Separated binary scales]\label{lem:905-energy}
Suppose the conclusion of Lemma~\ref{lem:905-curves} holds.  Every finite
$Y$-valued Walsh--Paley martingale $(M_k)_{k=0}^N$ with $M_0=0$ and
$\|M_N\|_\infty\leq1$ satisfies
\begin{equation}\label{eq:905-energy}
 \sum_{k=1}^N\mathbb E\|M_k-M_{k-1}\|^q\leq A,
 \qquad A=36\cdot12^qK.
\end{equation}
\end{lemma}
\begin{proof}
Represent its differences as
\[
 M_k(\epsilon_1,\ldots,\epsilon_k)
 =M_{k-1}(\epsilon_1,\ldots,\epsilon_{k-1})
 +\epsilon_k a_k(\epsilon_1,\ldots,\epsilon_{k-1}).
\]
All values of every $M_k$ belong to $B_Y$.  Choose positive integers
$m_1<\cdots<m_N$, and put $h_k=2^{-m_k}$.  Require
\[
 \frac32h_1\leq b,\qquad h_{k+1}\leq\frac14h_k,
 \qquad
 h_{k+1}\leq\frac{h_k\delta_k}{32}
 \quad\text{whenever }\delta_k\text{ is defined},
\]
where $\delta_k$ is the least positive norm among the finitely many
values of $a_k$.  If all its values vanish, omit the last requirement.
Such a finite sequence always exists.  Its gaps may depend on the
martingale, but none of our final constants do.

On $[0,1)$ let $r_m(t)=(-1)^{\lfloor2^mt\rfloor}$ and set
\[
 g_k(t)=M_k(r_{m_1}(t),\ldots,r_{m_k}(t)),\qquad
 F_k(t)=\int_0^t g_k(v)\,dv\quad(0\leq t\leq1).
\]
Extend $F_k$ by zero off $[0,1]$.  The selected binary digits are
independent fair signs, so $\int_0^1g_k=\mathbb EM_k=0$.
Thus these extensions are continuous, $1$-Lipschitz, supported in
$[0,1]$, and have supremum norm at most one.  Write $F=F_N$.

For $k<N$, on each dyadic cell of length $2h_{k+1}$ the earlier selected
digits are fixed and conditional averaging cancels the remaining
martingale differences.  Hence the integral of $g_N-g_k$ over every
such cell is zero.  Since its norm is at most two,
\begin{equation}\label{eq:905-tail}
 \|F-F_k\|_\infty\leq4h_{k+1}.
\end{equation}

Fix a full period $I=[2jh_k,(2j+2)h_k]$ of $r_{m_k}$ and write
$c=(2j+1)h_k$.  All earlier selected digits are constant on $I$.
There is an affine function $L_I$ such that
\[
 F_k(v)=L_I(v)+a_k(I)(h_k-|v-c|),\qquad v\in I.
\]
Let $\psi$ equal $+1$ on the middle half of $I$, and $-1$ on its two
outer quarters.  It has zero zeroth and first moments, and
\[
 \frac1{|I|}\int_I\psi(v)(h_k-|v-c|)\,dv=\frac{h_k}{4}.
\]
For every affine $\Lambda:\mathbb R\to Y$, the triangle and H\"older
inequalities imply
\[
 \left(\frac1{|I|}\int_I\|F_k-\Lambda\|^q\right)^{1/q}
 \geq\left\|\frac1{|I|}\int_I\psi(F_k-\Lambda)\right\|
 =\frac{h_k\|a_k(I)\|}{4}.
\]
This argument allows affine approximants with values outside the span
of the martingale.  Using \eqref{eq:905-tail} and the choice of
$h_{k+1}$ gives
\begin{equation}\label{eq:905-tent}
 \inf_\Lambda\left(\frac1{|I|}\int_I\|F-\Lambda\|^q\right)^{1/q}
 \geq\frac{h_k\|a_k(I)\|}{8}.
\end{equation}
For $k=N$ the tail is zero, and when $a_k(I)=0$ the bound is immediate.

If $|s-c|\leq h_k/4$ and $5h_k/4\leq u\leq3h_k/2$, the interval
$[s-u,s+u]$ contains $I$ and has length at most $3h_k$.  Thus
\[
 \beta_F(s,u)^q\geq\frac23
 \left(\frac{h_k\|a_k(I)\|}{8}\right)^q.
\]
The center windows for a fixed scale are disjoint, the scale windows
for distinct $k$ are disjoint, and all lie below $b$.  Each period
therefore contributes at least
\[
 \frac23\left(\frac{h_k\|a_k(I)\|}{8}\right)^q
 \frac{h_k}{2}\frac{h_k}{4}
 \left(\frac{3h_k}{2}\right)^{-q-1}
 =\frac{h_k\|a_k(I)\|^q}{18\cdot12^q}
\]
to $\mathcal D_b(F)$.  The preceding binary digits have the uniform
joint distribution over these periods, so
\[
 \sum_{I\text{ at scale }k}2h_k\|a_k(I)\|^q
 =\mathbb E\|M_k-M_{k-1}\|^q.
\]
Summing the disjoint contributions and applying
Lemma~\ref{lem:905-curves} proves \eqref{eq:905-energy}.
\end{proof}

\begin{lemma}[Renorming at the same exponent]\label{lem:905-renorm}
If \eqref{eq:905-energy} holds, then $Y$ admits an equivalent uniformly
convex norm of power type $q$.
\end{lemma}
\begin{proof}
For a unit-ball-valued martingale starting at any deterministic
$x\in B_Y$, subtracting $x$ and dividing by two permits the use of
\eqref{eq:905-energy}.  Its original energy is therefore at most
$2^qA$.  Put $H=\max\{1,2^qA\}$ and define
\[
 \Phi(x)=\sup\left\{\mathbb E\sum_{k=1}^N
 \|M_k-M_{k-1}\|^q:
 M_0=x,\ (M_k)\text{ a finite Walsh--Paley martingale in }B_Y\right\}.
\]
The zero-step martingale is allowed.  We have $0\leq\Phi\leq H$ and
$\Phi(-x)=\Phi(x)$.  If $Y$ is complex, the same argument gives
$\Phi(\lambda x)=\Phi(x)$ whenever $|\lambda|=1$, so the norm
constructed below is also a complex norm.  Join arbitrarily near-optimal martingales starting
at $x,y$ below a first fair split from $(x+y)/2$ to $x$ and $y$.
Their horizons can be equalized by constant steps.  Taking suprema gives
\begin{equation}\label{eq:905-bellman}
 \Phi\left(\frac{x+y}{2}\right)
 \geq\frac{\Phi(x)+\Phi(y)}2+\frac{\|x-y\|^q}{2^q}.
\end{equation}
Let $u(x)=\Phi(0)-\Phi(x)$.  Then $u$ is even, $u(0)=0$, and
\[
 0\leq u(x)\leq H,\qquad u(x)\geq\|x\|^q,
\]
where both lower bounds follow from \eqref{eq:905-bellman} at $x,-x$.
Moreover
\begin{equation}\label{eq:905-strong-convexity}
 u\left(\frac{x+y}{2}\right)
 \leq\frac{u(x)+u(y)}2-\frac{\|x-y\|^q}{2^q}.
\end{equation}
A locally bounded midpoint-convex function is continuous and convex
on the open unit ball.  Convexity on radial segments, first for dyadic
fractions and then by continuity, gives $u(x)\leq H\|x\|$ there.

Set $r=2^{-1/q}$, $a=1/(2H)$, and
\[
 K=\{x\in B_Y:u(x)\leq1/2\}.
\]
The estimates just proved imply $aB_Y\subseteq K\subseteq rB_Y$.
Since $r<1$, continuity in the open ball implies that $K$ is a closed,
symmetric convex body.  Its Minkowski functional $p$ is an equivalent
norm satisfying
\begin{equation}\label{eq:905-norm-comparison}
 a p(z)\leq\|z\|\leq r p(z).
\end{equation}
For completeness, we extract a power bound for its modulus.
Write $r'=(1+r)/2$ and $L=H/(1-r')=2H/(1-r)$.
The convex function $u$, with values in $[0,H]$ on $B_Y$, is
$L$-Lipschitz on $r'B_Y$.  Indeed, extend a line segment from either
endpoint a distance arbitrarily close to $1-r'$ within the unit ball,
and compare secant slopes using its range bound.

For $x,y\in K$ put $m=(x+y)/2$ and
$\Delta=(\|x-y\|/2)^q\leq r^q=1/2$.
By \eqref{eq:905-strong-convexity}, $u(m)\leq1/2-\Delta$.
Let $\eta=\Delta/(4L)$ and $z=m/(1-\eta)$.  Since
$\eta\leq1/(8L)$,
\[
 \|z\|\leq r+2\eta\leq r+\frac1{4L}<r',\qquad
 \|z-m\|\leq2\eta.
\]
It follows that $u(z)\leq1/2-\Delta+2L\eta\leq1/2$, so $z\in K$.
Consequently $p(m)\leq1-\eta$.  Whenever $p(x-y)\geq\varepsilon$,
\eqref{eq:905-norm-comparison} implies
\[
 p\left(\frac{x+y}{2}\right)
 \leq1-\frac{a^q}{2^{q+2}L}\varepsilon^q.
\]
This proves the asserted modulus of convexity at exponent $q$.
\end{proof}

\subsection{Passage to arbitrary probability spaces at the same exponent}

The final implication is the classical martingale renorming theorem
\cite{Pisier1975}, also recalled in \cite[Theorem~1]{HM2019}.
We include its needed direction directly to keep the exponent and the
probability-space quantifiers explicit.

\begin{lemma}\label{lem:905-martingale}
Suppose a norm $p$ satisfies $\delta_p(\varepsilon)\ge c\varepsilon^q$
for $0\le\varepsilon\le2$, where $q\ge2$ and $c>0$. Set
\[
 \kappa=\min\left\{4^{-q},\frac{c}{2\cdot3^q}\right\}.
\]
For every finite $L^q$ martingale on any probability space,
\[
 2\kappa\sum_{k=1}^{N-1}\mathbb E\,p(M_{k+1}-M_k)^q
 \le \mathbb E\,p(M_N)^q.
\]
\end{lemma}

\begin{proof}
We first prove the homogeneous convexity inequality
\begin{equation}\label{eq:905-power-midpoint}
 \frac{p(x)^q+p(y)^q}{2}
 -p\left(\frac{x+y}{2}\right)^q
 \ge\kappa p(x-y)^q.
\end{equation}
Write $A=p(x)\ge B=p(y)$, $d=p(x-y)$, and $t=(A+B)/2$.
If $A-B\ge d/2$, the scalar Clarkson inequality gives
\[
 \frac{A^q+B^q}{2}-t^q
 \ge\left(\frac{A-B}{2}\right)^q\ge4^{-q}d^q,
\]
and $p((x+y)/2)\le t$. This proves \eqref{eq:905-power-midpoint}
in that case. For completeness, the scalar inequality follows by
writing $A=t+s$, $B=t-s$: for $z\ge1$ the function
$(z+1)^q+(z-1)^q-2z^q$ is increasing, and at $z=1$ it is
$2^q-2\ge2$.

If $A-B<d/2$, the case $d=0$ is trivial, and otherwise $B>t/2$,
$A<3t/2$. Let $u=x/A$ and $v=y/B$. The triangle inequality implies
\[
 p(u-v)\ge\frac{d-(A-B)}A\ge\frac{d}{3t}.
\]
Use the modulus bound and
$(x+y)/2=(A-B)u/2+B(u+v)/2$ to obtain
\[
 p\left(\frac{x+y}{2}\right)
 \le t-Bc\left(\frac{d}{3t}\right)^q
 \le t-\frac{c\,d^q}{2\cdot3^q t^{q-1}}.
\]
For $0\le s\le t$, $t^q-s^q\ge t^{q-1}(t-s)$. Together with
$(A^q+B^q)/2\ge t^q$, this proves
\eqref{eq:905-power-midpoint} in the second case.

Put $F(x)=p(x)^q$. Regard $Y$ as a real space when taking supporting
functionals (use the real part of a complex norming functional). Choose a subgradient $\ell_x$ of $F$ at $x$.
Such a functional is supplied by a norming functional for $x$:
for $x\ne0$, take $\ell_x=q p(x)^{q-1}x^*$ with
$p^*(x^*)=1$ and $x^*(x)=p(x)$, and put $\ell_0=0$.
The supporting-functional inequality at $x$, evaluated at $(x+y)/2$
and combined with
\eqref{eq:905-power-midpoint}, gives
\begin{equation}\label{eq:905-tangent}
 F(y)\ge F(x)+\ell_x(y-x)+2\kappa p(y-x)^q.
\end{equation}
For a finite-valued martingale, choose $\ell_{M_k}$ separately at its
finitely many values. Conditional expectation of the linear term in
\eqref{eq:905-tangent} is zero. Summing over $k$ gives
\[
 2\kappa\sum_{k=1}^{N-1}\mathbb E\,p(M_{k+1}-M_k)^q
 \le\mathbb E F(M_N)-\mathbb E F(M_1)
 \le\mathbb E F(M_N).
\]

Finally, finite-valued martingales suffice even on an arbitrary probability
space. For each $k$ choose a finite-valued $\mathcal F_k$-measurable
function $s_k$ approximating $M_k$ in $L^q$ within $\epsilon$.
Let $\mathcal A_k=\sigma(s_1,\ldots,s_k)\subseteq\mathcal F_k$ and set
$N_k=\mathbb E[M_N\mid\mathcal A_k]$. This is a finite-valued martingale,
and the tower identity gives
$N_k=\mathbb E[M_k\mid\mathcal A_k]$. Since $s_k$ is
$\mathcal A_k$-measurable, conditional-expectation contraction yields
\[
 \|N_k-M_k\|_{L^q(p)}\le2\epsilon.
\]
The finite-valued inequality therefore passes to the original finite
martingale as $\epsilon\to0$. These conditional expectations are over
finite sigma-algebras and require no extra property of the Banach space.
\end{proof}

\begin{proof}[Proof of Theorem~\ref{thm:905}]
Apply the first three lemmas to obtain the equivalent norm $p$, with
$a p(z)\le\|z\|\le r p(z)$ and
$c=a^q/(2^{q+2}L)$. Lemma~\ref{lem:905-martingale} then gives
\[
 \left(\sum_{k=1}^{N-1}\mathbb E\|M_{k+1}-M_k\|^q\right)^{1/q}
 \le\frac{r}{a}(2\kappa)^{-1/q}
       \bigl(\mathbb E\|M_N\|^q\bigr)^{1/q}.
\]
This is the required martingale cotype inequality on every probability
space, with precisely the original exponent $q$.
\end{proof}

The decisive feature is that scale gaps may depend on each finite
martingale.  This makes every later primitive uniformly smaller than
the affine-error signal of each earlier difference, while the integral
with measure $du/u^{q+1}$ assigns the same energy weight at every chosen
scale.  The Bellman construction retains the exact exponent, so no
passage to $q+\varepsilon$ occurs.


\begin{thebibliography}{99}
\small
\setlength{\itemsep}{3pt}

\bibitem{HN2019}
T.~Hyt\"onen and A.~Naor,
\emph{Heat flow and quantitative differentiation},
J. Eur. Math. Soc. \textbf{21} (2019), 3415--3466.
\href{https://doi.org/10.4171/JEMS/906}{doi:10.4171/JEMS/906}.
Questions 10 and 19, and Theorem 9.

\bibitem{HM2019}
T.~Hyt\"onen and J.~Merikoski,
\emph{Vector-Valued Local Approximation Spaces},
J. Fourier Anal. Appl. \textbf{25} (2019), 299--320.
\href{https://doi.org/10.1007/s00041-018-9598-2}{doi:10.1007/s00041-018-9598-2}.

\bibitem{Pisier1975}
G.~Pisier,
\emph{Martingales with values in uniformly convex spaces},
Israel J. Math. \textbf{20} (1975), 326--350.
\href{https://doi.org/10.1007/BF02760337}{doi:10.1007/BF02760337}.
\end{thebibliography}
\end{document}
